\documentclass[smaller]{beamer}
\usepackage{booktabs}
\usepackage{xcolor}
\usepackage{multirow}

% --- THEME SETTINGS ---
\usetheme{Madrid}
\setbeamertemplate{navigation symbols}{}
\setbeamertemplate{footline}[frame number]

% --- PRESENTATION INFO ---
\title{Search Ad Optimization --- Research Update}
\date{February 19, 2026}

\begin{document}

\begin{frame}
\titlepage
\end{frame}

% ============================================================================
% Implementation Progress
% ============================================================================
\begin{frame}{Implementation Progress}
\framesubtitle{Google Ads API Integration}

Currently focusing on automating the inputs and outputs through the Google Ads API:

\medskip
\textbf{Inputs} (downloading data):
\begin{itemize}
    \item Keyword-level daily performance (clicks, costs).
    \item Campaign-level performance (clicks, purchases, by segment).
    \item Historical keyword metrics (search volume, competition).
    \item Previously downloaded manually; now being automated for daily retrieval.
\end{itemize}

\medskip
\textbf{Outputs} (uploading solution):
\begin{itemize}
    \item Campaign daily budgets (by region $\times$ match type).
    \item Base bids for each keyword--match type--region combination.
    \item Bid adjustments (hour of day, device, age, location).
\end{itemize}

\end{frame}

% ============================================================================
% LLM Embedding Option — What It Is
% ============================================================================
\begin{frame}{LLM Keyword Scoring}
\framesubtitle{Alternative to BERT embeddings}

\textbf{Motivation:} Replace generic BERT embeddings (768-d $\to$ 50-d via SVD) with a domain-specific relevance score from an LLM.

\medskip
\textbf{Approach:}
\begin{itemize}
    \item Use Qwen3-8B to score each keyword on a 1--5 relevance scale.
    \item The score replaces the 50 embedding dimensions with a single feature.
    \item Scoring is based on a rubric that accounts for purchase intent, learning intent, and irrelevance signals.
\end{itemize}

\medskip
\textbf{Prompt structure} (keyword = ``machine learning course'', course = ML):
\begin{enumerate}
    \item Start with Base Score = 3.
    \item $-1$ if keyword contains ``free'', ``cheap'', ``youtube'', ``login'', or is unrelated.
    \item $-1$ if purely informational (``what is'', definitions, news).
    \item $+1$ if implies learning intent (``course'', ``training'', ``tutorial'').
    \item $+1$ if implies high purchase intent (``certification'', ``bootcamp'', ``university'').
\end{enumerate}
Apply all modifiers cumulatively. Score clamped to $[1, 5]$.

\end{frame}

% ============================================================================
% Comparison: BERT vs LLM
% ============================================================================
\begin{frame}{BERT vs.\ LLM Scores}
\framesubtitle{Gen AI course, Dec 1--31 backtest}

\begin{table}[htbp]
\centering
\resizebox{\textwidth}{!}{%
\begin{tabular}{lcccccccc}
\toprule
 & \multicolumn{5}{c}{Optimized (daily avg)} & \multicolumn{3}{c}{Improvement (vs.\ Baseline)} \\
\cmidrule(lr){2-6} \cmidrule(lr){7-9}
Experiment & Clicks & Purch & Cost & Clicks/\$ & Kws & Clicks & Purch & Clicks/\$ \\
\midrule
\textbf{BERT} & \textbf{242.5 $\pm$ 11.5} & \textbf{0.16 $\pm$ 0.01} & 362.83 & \textbf{0.668} & 101 & \textbf{+31.2\%} & \textbf{+42.3\%} & \textbf{+25.8\%} \\
LLM & 176.8 $\pm$ 7.5 & 0.11 $\pm$ 0.00 & 362.84 & 0.487 & 53 & +6.0\% & +7.2\% & +1.7\% \\
\bottomrule
\end{tabular}
}
\end{table}

\medskip
\textbf{Key takeaway:} BERT embeddings substantially outperform LLM relevance scores in this backtest setting.

\end{frame}

% ============================================================================
% Model Fit Comparison
% ============================================================================
\begin{frame}{BERT vs.\ LLM Scores}
\framesubtitle{Model quality comparison}

\begin{table}[htbp]
\centering
\begin{tabular}{lccc}
\toprule
Experiment & MSE & $R^2$ \\
\midrule
\textbf{BERT} & \textbf{4.40 $\pm$ 0.04} & \textbf{0.927 $\pm$ 0.000} \\
LLM & 5.44 $\pm$ 0.11 & 0.913 $\pm$ 0.001 \\
\bottomrule
\multicolumn{3}{l}{\scriptsize Averaged over 31 backtest days (Dec 1--31).} \\
\end{tabular}
\end{table}

\medskip
\textbf{Discussion:}
\begin{itemize}
    \item BERT provides 50 embedding dimensions capturing semantic variation among keywords; the LLM approach compresses this to a single 1--5 score.
    \item The LLM score is more interpretable but the existing set of keywords all have similar scores.
    \item Could be useful in the future when we have results for a broader set of keywords with varying LLM scores.
\end{itemize}

\end{frame}

% ============================================================================
% Next Steps
% ============================================================================
\begin{frame}{Next Steps}

\begin{enumerate}
    \item \textbf{Continue implementation:} Google Ads API integration for automated daily data download and bid upload.

    \item \textbf{Fine-tune BERT embedding dimensions:} Currently using 50 SVD components; experiment with different values (e.g., 20, 100, 200) to find the optimal dimensionality.

    \item \textbf{Add robustness via ensemble models:} Train multiple XGBoost models (with different bootstrap samples) and average their predictions, as described in the reference paper.

    \item \textbf{Validate bid multiplier:} The $\times 1.3$ second-price auction adjustment needs empirical validation once live bidding begins.
\end{enumerate}

\end{frame}

\end{document}
